\uuid{8QU1}
\exo7id{2329}
\titre{exo7 2329}
\auteur{matexo1}
\organisation{exo7}
\datecreate{2002-02-01}
\isIndication{false}
\isCorrection{false}
\chapitre{Calcul d'intégrales}
\sousChapitre{Intégration par parties}
\cpge{mpsi-mp2i.primitives-equations-differentielles.a, pcsi.techniques-calcul-primitives.a, ptsi.techniques-calcul-primitives.a, mpsi-mp2i.limites-continuite.a, pcsi.limites-continuite.a, ptsi.limites-continuite.a, mpsi-mp2i.integration.c, pcsi.integration.b, ptsi.integration.b, bcpst1.derivees-primitives-integrales.d, bcpst1.integration-segment.d}
\module{Analyse}
\niveau{L1}
\difficulte{}

\contenu{
\question{
Soit $x>0$ un r\'eel. Calculer les valeurs de
$$I(x) = \int_0^x{\arctan t \over 1+t^2} dt \qquad \mbox{\rm et}
\qquad J(x) = \int_0^x{\arctan t \over (1+t)^2} dt. $$
Quelles sont leurs limites quand $x \to  +\infty$\,?
}
}
